\documentclass[12pt,a4paper]{article}

\usepackage{cmap}
\usepackage[T2A]{fontenc}
\usepackage[utf8]{inputenc}
\usepackage[russian]{babel}
\usepackage{amsmath,amssymb}
\usepackage[a4paper,margin=2cm]{geometry}

\setlength{\parindent}{0pt}
\setlength{\parskip}{0.45em}
\sloppy

\begin{document}

\begin{center}
{\Large\bfseries Билет 4: вопросы для проверки знаний}
\end{center}

\noindent\fbox{%
  \begin{minipage}{\dimexpr\linewidth-2\fboxsep-2\fboxrule\relax}
  \normalfont
Частные производные высших порядков. Независимость смешанной частной
производной от порядка дифференцирования. Дифференциалы высших порядков.
Формула Тейлора для функций нескольких переменных с остаточным членом в формах
Лагранжа и Пеано.
  \end{minipage}%
}

\section*{Вопросы на понимание}

\begin{enumerate}
\item Чем смешанная производная второго порядка отличается по смыслу от повторной производной по одной и той же переменной?
\item Почему существование обеих смешанных производных в точке само по себе еще не гарантирует их равенства?
\item Какую роль играет непрерывность смешанных производных в теореме о независимости от порядка дифференцирования?
\item Зачем в формулировках о смешанных производных обычно требуют существование производных в окрестности точки, а не только в самой точке?
\item Может ли порядок дифференцирования стать важным, если у функции есть частные производные нужного порядка, но нет нужной непрерывности? Объясните идею на уровне контрпримера.
\item Почему при вычислении $d^k f$ величины $dx_1,\dots,dx_n$ считают постоянными?
\item Откуда в формуле для $d^k f(x,y)$ появляются биномиальные коэффициенты, если смешанные производные одного порядка равны?
\item Почему в многочлене Тейлора для функции нескольких переменных в дифференциалах полагают $dx_i=\Delta x_i$?
\item Как переход к функции одной переменной $\varphi(t)=f(x^0+t\Delta x)$ помогает получить формулу Тейлора для функции нескольких переменных?
\item Почему для остатка в форме Лагранжа важна точка $x^0+\theta\Delta x$ на отрезке между $x^0$ и $x$?
\item Чем остаток в форме Пеано концептуально отличается от остатка в форме Лагранжа?
\item Почему для формы Пеано естественно требовать непрерывность производных до порядка $m$, а для формы Лагранжа появляется производная порядка $m+1$?
\end{enumerate}

\section*{Источники для сверки}

\texttt{target/answer\_question\_04.tex};
\texttt{IvGE\_1.pdf}, стр. 191--199.

\end{document}
